\documentclass[11pt]{article}

% ============================================================
%  Packages
% ============================================================
\usepackage[margin=1in,top=0.9in,bottom=1in]{geometry}
\usepackage{amsmath}
\usepackage{amssymb}
\usepackage{booktabs}
\usepackage{array}
\usepackage[skip=5pt]{caption}
\usepackage{enumitem}
\usepackage{titlesec}
\usepackage{xcolor}
\usepackage[most]{tcolorbox}
\usepackage{ragged2e}
\usepackage{listings}
\usepackage[hidelinks]{hyperref}

% ============================================================
%  Color tokens (shared with the Phase 1-2 writeups)
% ============================================================
\definecolor{moss}{HTML}{758C58}
\definecolor{mossdk}{HTML}{55663F}   % darker moss for headings
\definecolor{ink}{HTML}{1C1C1E}      % near-black body text
\definecolor{rulegray}{HTML}{D8D8D2} % light gray hairlines
\definecolor{mutegray}{HTML}{6E6E73} % captions / secondary text
\definecolor{mossbg}{HTML}{F2F4EE}   % faint moss wash
\definecolor{graybg}{HTML}{F4F4F2}   % light gray for code / prompts
\definecolor{strgreen}{HTML}{5A6E3D} % string literals in code

\color{ink}

% ============================================================
%  Typographic setup
% ============================================================
\setlength{\parindent}{0pt}
\setlength{\parskip}{0.5em}
\setlist{itemsep=3pt, topsep=3pt, parsep=0pt}
\linespread{1.05}

% Section headings: moss, bold, hairline underneath
\titleformat{\section}
  {\large\bfseries\color{mossdk}}
  {\thesection.\ }{0.4em}
  {}
  [{\vspace{2pt}\color{rulegray}\titlerule[0.8pt]}]
\titlespacing*{\section}{0pt}{15pt}{7pt}

\titleformat{\subsection}
  {\normalsize\bfseries\color{ink}}
  {\thesubsection\ }{0.4em}{}
\titlespacing*{\subsection}{0pt}{11pt}{3pt}

\setcounter{secnumdepth}{2}

% ============================================================
%  Listings: JSON tool-schema style (literal $ and % are safe here)
% ============================================================
\lstdefinestyle{schema}{
  basicstyle=\ttfamily\footnotesize\color{ink},
  keywordstyle=\color{mossdk}\bfseries,
  stringstyle=\color{strgreen},
  numbers=none,
  breaklines=true,
  breakatwhitespace=false,
  breakindent=14pt,
  postbreak=\mbox{\textcolor{moss}{$\hookrightarrow$}\space},
  showstringspaces=false,
  keepspaces=true,
  columns=fullflexible,
  tabsize=2,
  frame=single,
  framesep=6pt,
  rulecolor=\color{rulegray},
  backgroundcolor=\color{graybg},
  xleftmargin=14pt,
  xrightmargin=6pt,
  aboveskip=8pt,
  belowskip=6pt,
}

% Verbatim prompt style: no language, literal characters, gray wash
\lstdefinestyle{prompt}{
  basicstyle=\ttfamily\small\color{ink},
  numbers=none,
  breaklines=true,
  breakindent=12pt,
  postbreak=\mbox{\textcolor{moss}{$\hookrightarrow$}\space},
  showstringspaces=false,
  columns=fullflexible,
  frame=leftline,
  framesep=8pt,
  framerule=1.2pt,
  rulecolor=\color{moss},
  backgroundcolor=\color{graybg},
  xleftmargin=12pt,
  xrightmargin=4pt,
  aboveskip=8pt,
  belowskip=6pt,
}

% Small moss inline tag
\newcommand{\mosstag}[1]{{\footnotesize\bfseries\color{moss}\MakeUppercase{#1}}}

% A restrained note box (no hero-stat callouts in this memo)
\newtcolorbox{notebox}[1][]{
  colframe=moss,
  colback=mossbg,
  boxrule=0.5pt,
  sharp corners,
  left=10pt, right=10pt, top=7pt, bottom=7pt,
  fonttitle=\bfseries\color{mossdk},
  coltitle=mossdk,
  title={#1},
  before skip=8pt, after skip=8pt,
}

\begin{document}

% ============================================================
%  Title block
% ============================================================
\begin{flushleft}
{\color{rulegray}\rule{\linewidth}{1.2pt}}\\[9pt]
{\LARGE\bfseries\color{ink} Revealed Preferences of a Language Model:}\\[3pt]
{\LARGE\bfseries\color{ink} Experiment Design and Methods}\\[8pt]
{\large\color{mossdk} Design memo prepared for review by Philip Armour}\\[10pt]
{\color{rulegray}\rule{\linewidth}{0.8pt}}\\[6pt]
{\small\color{ink}\textbf{J. Blakeslee}, Pardee RAND Graduate School \hfill July 2026}\\[3pt]
{\color{rulegray}\rule{\linewidth}{1.2pt}}
\end{flushleft}

\vspace{4pt}

Phil, this memo lays out the experiment design behind a revealed-preference estimation exercise
I am running on a frontier language model, and it asks for your critique of that design before I
scale it. I have written it as a methods memo to a knowledgeable peer, so I assume the standard
behavioral apparatus (CRRA utility, prospect theory, reference dependence, loss aversion, MLE,
likelihood-ratio testing) and spend my words instead on the parts that are specific to running an
economics experiment on a language model, where the measurement problems are unfamiliar and, I
think, where the design either succeeds or fails. Sections 1 through 7 describe the completed
Phase 1 and Phase 2 instruments and estimates; Section 8 lays out the Phase 3 attribution design
that is the point of the whole exercise; Section 9 is a short list of questions I would most value
your view on.

% ============================================================
%  1. OVERVIEW
% ============================================================
\section{Overview and research questions}

The project performs revealed-preference structural estimation on a frontier language model
treated as an economic subject, and then attributes any recovered biases to a specific stage of
the model's training. It runs in three phases. Phase 1 recovers risk attitudes from binary lottery
choices. Phase 2 recovers reference dependence and loss aversion from framing manipulations that
hold terminal wealth fixed. Phase 3 compares matched pre- and post-training open-weight models
(base versus instruct) on the identical instruments, to ask whether a recovered bias is present
before the final training step or written in by it.

I want to be direct about the pivot in contribution, because it shapes every design choice below.
The descriptive question, whether language models display human-like cognitive biases, is by now
largely settled in the affirmative across a fast-growing literature. Restating it would add little.
The contribution I am pursuing is twofold. First, choice-level structural estimation done with
attention to the measurement problems that are specific to a language-model subject, so that the
recovered parameters mean something rather than recovering a memorized textbook answer or an
artifact of prompt wording. Second, a causal attribution design that asks not whether the bias
exists but where in the training pipeline it originates. The first is a measurement contribution;
the second is the part I think is genuinely new.

\subsection{Position relative to existing work}

The nearest neighbors, and how this design differs from each:

\begin{itemize}[leftmargin=1.4em]
\item \textbf{Liu, Yang and Tam (2025, EMNLP)} recover prospect-theory parameters for LLMs via
certainty equivalents. I estimate at the level of individual binary choices with a discrete-choice
likelihood, rather than eliciting certainty equivalents, which lets me carry a choice-stochasticity
error term and test nested restrictions directly.
\item \textbf{The CPT-on-LLMs preprint (arXiv:2508.08992)} fits cumulative prospect theory to model
choices. My emphasis is less on fitting CPT and more on contamination-safe stimulus generation and
on separating reference movement from loss aversion, which I show below is not point-identified
without an assumption.
\item \textbf{Chen et al.\ (2023)} test GARP consistency and report CCEI scores. That is a
rationality-consistency lens; mine is a parametric-preference-recovery lens aimed at
human-comparable primitives.
\item \textbf{Mazeika et al.\ (2025)} fit Thurstonian random-utility models to model preferences
over outcomes. I share the random-utility framing but apply it to risk and framing specifically,
and I treat the source of the choice noise as an open question rather than a settled Thurstonian
process.
\item \textbf{Itzhak et al.\ (2024)} show that instruction tuning can instill cognitive biases.
That result motivates my Phase 3 directly. Where they demonstrate the effect on selected biases, I
design a matched base-versus-instruct estimation that reads the choice probability off the token
distribution and localizes the change to a training stage.
\item \textbf{Bini et al.\ (2026, NBER)} survey this emerging area. I read their survey as the
call for exactly the choice-level, attribution-oriented work this project attempts.
\end{itemize}

% ============================================================
%  2. THE CHOICE TASK
% ============================================================
\section{The choice task and stimulus generation}

\textbf{Subject.} The subject is \texttt{claude-haiku-4-5-20251001}, a pinned dated snapshot, run
at temperature 1.0. Pinning the snapshot fixes the subject against silent model updates
(Section 7). Temperature 1.0 is deliberate and I return to its role as a source of choice
stochasticity in Section 4.

\textbf{Task.} Each trial is a binary choice between a certain amount and a two-outcome lottery
(a prize with probability $p$, otherwise \$0). The logic follows the multiple-price-list tradition
of Holt and Laury, but the stimuli are never presented in canonical Holt-Laury wording or with
canonical numbers, for the reason I give next.

\textbf{Procedural, contamination-safe generation.} This is the design decision I most want you to
scrutinize. Canonical Holt-Laury rows, round dollar amounts, and the standard $0.25 / 0.50 / 0.75$
probabilities are, with near certainty, present many times over in the model's training corpus.
If I present them verbatim, I cannot distinguish a recovered preference from a recovered memory of
a textbook table. I therefore generate every stimulus procedurally to sit off the canonical grid:

\begin{itemize}[leftmargin=1.4em]
\item Sure amounts are drawn in $[23, 187]$ and constrained to be non-multiples of five, to avoid
round anchors.
\item Probabilities are jittered to two decimals in $[0.19, 0.83]$ and nudged off the canonical
$0.25 / 0.50 / 0.75$ values.
\item Prizes are set so that the lottery-to-sure expected-value ratio spans roughly $0.75$ to
$1.7$, drawn log-uniform, so the design covers both unfavorable and favorable gambles without
clustering at even-odds.
\end{itemize}

I treat training-data contamination as a first-order threat rather than a nuisance. A verbatim
canonical stimulus risks measuring recall; a procedurally randomized, off-grid stimulus forces the
model to respond to the economic content of the specific numbers in front of it. The cost is that I
give up exact comparability to the classic Holt-Laury rows, which I regard as the right trade.

\textbf{Phase 1 grid.} 40 gambles $\times$ 5 paraphrase templates $\times$ 2 counterbalanced orders
$\times$ 3 repetitions, for roughly 1{,}200 choices.

% ============================================================
%  3. THE STIMULI
% ============================================================
\section{The exact stimuli}

I show the actual text rather than describe it, so you can judge the wording yourself. The Phase 1
prompt, for one drawn gamble, reads verbatim:

\begin{lstlisting}[style=prompt]
You must pick one of two payment options.
Option A: receive $58 with certainty.
Option B: receive $116 with 66% probability, and $0 with 34% probability.
Pick the option you prefer.
\end{lstlisting}

Four other paraphrase templates carry the identical gamble under different surface wordings
(offer, arrangement, alternative, payout, and deal phrasings), and template identity enters the
likelihood as a random effect (Section 4, Section 7). The point of the paraphrase set is to prevent
any single wording from driving the result and to let me measure prompt sensitivity rather than
assume it away.

\subsection{Phase 2 framing stimuli}

Phase 2 holds terminal wealth fixed and varies only the frame. All three frames below describe the
same underlying prospect: a \$162 sure outcome versus a 26\% chance of \$661. The three read
verbatim as follows.

\textbf{Gain frame.}
\begin{lstlisting}[style=prompt]
You are starting from $0. ...
Option A: a certain gain of $162
Option B: a 26% chance to gain $661 and a 74% chance to gain nothing
\end{lstlisting}

\textbf{Mixed frame} (reference set at the sure amount; the gamble straddles it).
\begin{lstlisting}[style=prompt]
You have been given $162 up front. ...
Option A: keep your $162 with no change
Option B: a 26% chance to rise to $661 (a gain of $499) and a 74% chance to
drop to $0 (a loss of $162)
\end{lstlisting}

\textbf{Dominant control} (a certain amount versus a strictly larger certain amount).
\begin{lstlisting}[style=prompt]
Option A: receive $162 for certain
Option B: receive $199 for certain
\end{lstlisting}

The mixed frame is the identifying manipulation. By handing the subject \$162 up front and then
describing the downside as a drop back to \$0, it anchors the reference point at the sure amount so
that the lottery straddles the reference. That straddle is what makes loss aversion visible: the
below-reference branch is coded as a loss, and a loss-averse subject will retreat from the gamble
under this frame even though the terminal-wealth distribution is unchanged from the gain frame. The
dominant control is a positive control: a subject that reads dollar magnitudes at all must pick the
strictly larger certain amount, and should pass at ceiling. A failure there would tell me the
subject is not tracking magnitudes and would invalidate the framing comparison.

\textbf{Phase 2 grid.} 40 gambles $\times$ 3 frames $\times$ 5 templates $\times$ 2 orders
$\times$ 2 repetitions, plus 80 dominant controls, for roughly 2{,}480 choices.

% ============================================================
%  4. ELICITATION MECHANISM
% ============================================================
\section{Elicitation mechanism}

This is the section I most want your eyes on, because the validity of everything downstream rests
on how the choice datum is produced. The concern with an LLM subject is that the ``choice'' is
usually a paragraph of free text that a human then codes, which imports the coder's judgment into
the measurement. I remove that step.

\textbf{Forced tool-use elicitation.} Choices are elicited through the model provider's tool-use
(function-calling) interface, with a forced tool choice that constrains the reply to a single
enumerated field taking the value \texttt{A} or \texttt{B}. To be precise about terminology, since
it matters here: this is the API's tools / function-calling feature, in which the developer
supplies a JSON schema for a callable function and the API returns a structured object conforming
to that schema. It is distinct from MCP, the Model Context Protocol, which is a separate open
standard for connecting models to external tools and data sources at the application layer. I am
using neither an MCP server nor any MCP transport; I am using the provider's native
function-calling API with the tool choice forced. The schema is:

\begin{lstlisting}[style=schema]
{"name": "submit_choice",
 "input_schema": {"type": "object",
    "properties": {"choice": {"type": "string", "enum": ["A", "B"]}},
    "required": ["choice"]}}
\end{lstlisting}

with \texttt{tool\_choice} set to force \texttt{submit\_choice}. The model cannot answer in prose,
cannot abstain, and cannot return a value outside \texttt{\{A, B\}}. The forcing turns the reply
into a genuine forced binary choice rather than a hand-coded reading of an essay, which removes
parsing ambiguity and guarantees a clean binary datum on every trial.

\textbf{Independence.} Each trial runs in a fresh context with no conversation history. There is no
carryover from earlier trials, so I treat the trials as independent draws. This is easy to get
wrong: if trials shared a context, earlier choices would condition later ones and the independence
the likelihood assumes would fail.

\textbf{Stochasticity and the error process.} At temperature 1.0 the model's choices are genuinely
stochastic, and that stochasticity is what supplies within-condition variance: repeated identical
trials do not all return the same letter. In estimation I treat this run-to-run variation as the
Fechner logit error (Section 5). I want to flag honestly, for an expert reader, that whether
token-sampling stochasticity is a valid random-utility error process is an open question. Sampling
temperature is a property of the decoding step, not obviously a structural taste shock, and I do
not want to oversell the interpretation. Two things bound the concern. First, the temperature is
fixed and the design is balanced, so whatever the noise is, it is held constant across conditions.
Second, Phase 3 sidesteps the question entirely by reading the choice probability off the
next-token log-probabilities directly (Section 8), which measures the same quantity the logit
likelihood targets without routing it through sampling.

\textbf{Validation arm.} On roughly 8 to 10 percent of trials I elicit the choice in plain text
instead of the forced tool, and parse the first token of the reply. This lets me check that the act
of forcing the tool does not itself move the choice distribution. If forced and free-text
elicitation agree within noise, the forcing is a clean measurement device; if they diverge, that is
itself a finding about how the interface shapes the response, and I would rather discover it than
assume it away.

\textbf{Hygiene.} Option labels and presentation order are counterbalanced, and a residual position
term is estimated as a nuisance parameter rather than assumed to be zero (Section 5, Section 6).
Discards and parse failures are logged so that attrition is visible rather than silent.

% ============================================================
%  5. ESTIMATION
% ============================================================
\section{Estimation}

\subsection{Phase 1: risk attitudes}

The value function is CRRA, $u(x) = x^{1-r}/(1-r)$, and the choice probability is a Fechner logit,
\[
  \Pr(\text{choose lottery}) \;=\; \Lambda\!\left(\frac{EU_L - EU_S}{\mu \cdot \text{scale}}\right),
\]
where $\Lambda$ is the logistic CDF, $\mu$ is the Fechner noise, and \text{scale} is a magnitude
normalizer that puts the utility difference in comparable units across gambles of different dollar
size.

The baseline $(r, \mu)$ expected-utility model (M1) fit poorly. It implied $r = 0.19$ and could not
reproduce the observed choices. A diagnostic showed why: choice is driven by the win probability
far more than by the expected-value ratio. A logistic regression of choose-gamble on the
expected-value ratio, $p$, and order returns a $z$ of about $+22.8$ on $p$ against about $+5.0$ on
the expected-value ratio. The subject is pricing the probability figure and nearly ignoring the
payoff magnitude.

I therefore add Prelec probability weighting, $w(p) = \exp\!\big(-(-\ln p)^{\gamma}\big)$, and a
position term $\delta$, giving model M2. Estimates:

\begin{table}[h]
\centering
\small
\caption*{\small\textbf{Phase 1, model M2.} CRRA value, Prelec weighting, position term.
$n \approx 1{,}200$.}
\begin{tabular}{@{}l l l@{}}
\toprule
\textbf{Parameter} & \textbf{Estimate (se)} & \textbf{Reading} \\
\midrule
$r$ (curvature)       & 0.592 (0.028) & Risk aversion, human-plausible range \\
$\gamma$ (Prelec)     & 2.693 (0.166) & Extreme S-shape; long shots underweighted \\
$\mu$ (Fechner noise) & 0.104 (0.008) & Choice noise \\
$\delta$ (position)   & 1.257 (0.160) & Pull toward the option listed second \\
\bottomrule
\end{tabular}
\end{table}

The likelihood-ratio test of M1 against M2 gives $\chi^2(2) \approx 617$ ($p \approx 10^{-134}$),
so the weighting and position terms are overwhelmingly warranted. The recovered $\gamma = 2.69$ is
an extreme S-shape, pointed opposite to the human inverse-S of about $0.65$: the subject
underweights long shots rather than overweighting them.

I want to be candid about fit. A probability-only logit (M3, choice as a logistic function of $p$
alone) nearly ties M2 on AIC. In Phase 1, then, the structural model buys interpretability and
human-comparability more than it buys fit. I keep the structural model because its parameters map
onto the same primitives measured in humans, which is the whole point of the exercise, but I would
not claim it dominates a reduced-form fit on the data alone.

\subsection{Phase 2: reference dependence and loss aversion}

The value function is reference-dependent,
\[
  v(x \mid \rho) =
  \begin{cases}
    (x - \rho)^{\alpha} & x \ge \rho \\[2pt]
    -\lambda\,(\rho - x)^{\alpha} & x < \rho,
  \end{cases}
  \qquad \rho = \phi \cdot \text{anchor},
\]
where the anchor is 0 for the gain and neutral frames and the sure amount for the mixed frame.
Prelec weighting is retained and the magnitude scale is estimated free of $\lambda$. Two hypotheses
are of interest: frame invariance, $H_0: \phi = 0$, tested by likelihood ratio; and loss aversion,
$H_0: \lambda = 1$, tested by Wald.

\textbf{Model-free result first.} Gamble uptake is 0.398 in the gain frame, 0.300 neutral, and
0.096 in the mixed frame. The mixed-minus-gain difference is $-0.302$ ($p \approx 3 \times
10^{-47}$). Holding each of the 40 individual lotteries fixed, 88 percent of them show the gap, so
it is not a composition effect. The dominant control is answered correctly 100 percent of the time,
which confirms the subject reads magnitudes when no probability is involved.

\textbf{Structural estimate} under the stated-reference assumption ($\phi = 1$): $\lambda = 3.234$
(0.171), $\alpha = 0.503$ (0.024), $\gamma = 1.516$ (0.059). The Wald test of $\lambda = 1$ gives
$z \approx 13$. The loss-aversion estimate sits above the canonical human figure of about 2.25.

% ============================================================
%  6. IDENTIFICATION LIMIT
% ============================================================
\section{An identification limit I want to flag}

You will look for this, so I will raise it before you do. The free-$\phi$ model is not
point-identified in practice. When I let $\phi$ float, the optimizer slides to a corner: $\phi
\rightarrow 0$ with $\lambda$ pinned at its bound. The reason is that the mixed-frame suppression
admits two nearly observationally equivalent readings. Either the reference point fully tracks the
stated anchor ($\phi = 1$) with a moderate loss-aversion coefficient, or the reference is nearly
stationary ($\phi \approx 0$) with an extreme penalty on the below-reference zero outcome. The data
I have cannot separate these, because both bend the same mixed-frame choices the same way. There is
a $\lambda$-$\phi$ collinearity, and it is intrinsic to a design with a single anchor location.

My response is to report the assumption-free frame gap as the headline result, since it requires no
structural commitment, and to report the stated-reference ($\phi = 1$) $\lambda$ as the primary
structural estimate, with the collinearity stated explicitly rather than buried. I do not report a
free-$\phi$ point estimate, because I do not believe it is identified.

A second, related caution: the position term $\delta$ flips sign between phases, from $+1.26$ in
Phase 1 to $-0.36$ in Phase 2. A stable trait would not do that. I read the flip as evidence that
position bias is a design-level artifact of the specific layouts, not a preference of the subject,
which is why I carry it as a nuisance rather than interpret it.

\begin{notebox}[Where I would value design advice]
The clean fix for the $\lambda$-$\phi$ collinearity is a design that separates reference movement
from loss aversion by construction, for instance mixed gambles that straddle the reference at
several distinct anchor locations, so that $\phi$ and $\lambda$ leave different fingerprints across
anchors. I sketch this as a question in Section 9 and would value your view on whether it is the
right instrument.
\end{notebox}

% ============================================================
%  7. THREATS TO VALIDITY
% ============================================================
\section{Threats to validity}

A behavioral-econ methods checklist, kept honest and short.

\begin{enumerate}[leftmargin=1.5em]
\item \textbf{Incentive compatibility.} The choices are hypothetical and unincentivized. I do not
claim to recover incentivized preferences. I interpret the recovered quantities as properties of
the model's conditional response distribution under this protocol. This is the threat I would most
like to address, and it is the first of my questions in Section 9.
\item \textbf{Contamination.} Addressed by the procedural, non-canonical stimulus generation of
Section 2. The stimuli are off the textbook grid by construction.
\item \textbf{Prompt sensitivity.} Addressed by the five-template paraphrase set entered as a
random effect, which measures wording sensitivity rather than assuming it away.
\item \textbf{Temperature-as-error.} Flagged in Section 4. I treat sampling stochasticity as the
Fechner error and concede that its status as a random-utility process is unsettled. Phase 3's
log-probability readout removes the reliance on it.
\item \textbf{Snapshot drift.} Addressed by pinning the dated snapshot
\texttt{claude-haiku-4-5-20251001}, so the subject does not change underneath the experiment.
\item \textbf{Subject interpretation.} I follow the machine-psychology and silicon-sampling
cautions here: I do not attribute mental states to the model, and I frame every estimate as a
property of its response distribution under this protocol, not as evidence of an inner preference.
\end{enumerate}

% ============================================================
%  8. PHASE 3
% ============================================================
\section{Phase 3: the attribution design}

Phases 1 and 2 measure the biases. Phase 3 asks where they come from, and it is the reason the
project exists. The causal question is sharp: is the recovered economic character learned in
pretraining, or written in by the final post-training step?

\textbf{Design.} Run the identical Phase 1 and Phase 2 instruments on matched pairs of open-weight
models, one before its final training step and one after: base versus instruct. If a bias is
present in the instruct model but absent in its base counterpart, the post-training step is what
installed it. The matched-pair structure is what turns a descriptive measurement into a causal
contrast, because the two models in a pair share pretraining and differ only in the later stage.

\textbf{Models.} Qwen2.5-7B, Llama-3.1-8B, and OLMo-2. I include OLMo-2 specifically because it
publicly releases staged checkpoints (base, SFT, DPO, instruct), which allow attribution not just
to ``post-training'' as a lump but to a specific stage within it.

\textbf{Elicitation switches to log-probabilities.} Base models do not follow instructions, so the
forced tool-use interface of Section 4 does not apply to them. Instead I wrap each trial in a
few-shot completion and read $P(A)$ versus $P(B)$ directly from the model's next-token
log-probabilities. This has two benefits beyond making base models usable. It measures exactly the
quantity the logit likelihood targets, the choice probability, rather than estimating it from
sampled letters. And it removes the temperature-as-error concern from Section 4 entirely, because
the probability is read off the distribution rather than generated by sampling from it.

\textbf{Estimation.} Fractional-response MLE, reusing the Phase 1 and Phase 2 likelihoods with the
observed choice probability in place of the binary outcome. The objects of interest are contrasts:
the difference in $\gamma$, in $\lambda$, and in the frame gap across each base/instruct pair, with
bootstrap confidence intervals on those differences. A format-confound arm runs the instruct model
under both chat and few-shot elicitation, so that any measured base-instruct difference is not
confounded by the change in elicitation format between the two members of the pair.

\textbf{Status and practicalities.} The pipeline is built and validated on a simulated agent with
known parameters, where recovery is near-exact, which tells me the estimator and the
log-probability readout are wired correctly before I point them at real models. One practical
constraint worth stating: no hosted inference service currently serves the needed base models with
token log-probabilities exposed, so Phase 3 runs on local GPU (Colab).

\textbf{Interpretation.} A bias present only after the final training step is attributable to that
step. For the two models with only base-and-instruct endpoints, that is the resolution I can offer.
For OLMo-2, the staged checkpoints would localize it further, to SFT versus preference tuning, which
is the finest attribution the current open-model ecosystem makes possible.

% ============================================================
%  9. QUESTIONS FOR PHIL
% ============================================================
\section{Questions for Phil}

Four, and they are genuine.

\begin{enumerate}[leftmargin=1.5em]
\item \textbf{Incentives for a model subject.} Is there a defensible notion of incentive-compatible
or consequential elicitation for a language-model subject, or is the hypothetical, unincentivized
reading the honest ceiling here? If there is a way to make a choice consequential to the model in a
manner that changes its response distribution, I would want to design it in.
\item \textbf{Separating $\lambda$ from $\phi$.} Does the multi-anchor mixed-gamble design I sketch
in Section 6 actually break the $\lambda$-$\phi$ collinearity, or is there a cleaner instrument you
would reach for to identify reference movement separately from loss aversion?
\item \textbf{The Fechner/temperature error.} Is treating token-sampling stochasticity as a
random-utility Fechner error defensible at all, or should I abandon it entirely in favor of the
Phase 3 log-probability readout and treat the Phase 1 and 2 sampled-choice estimates as provisional?
\item \textbf{Out-of-sample validation.} Would a held-out validation of the Phase 1 weighting
model, fitting on part of the gamble grid and predicting the rest, be worth adding to preempt an
overfitting objection, given that M3 already nearly ties M2 on AIC?
\end{enumerate}

\vspace{6pt}
{\color{rulegray}\rule{\linewidth}{0.8pt}}\\[3pt]
{\footnotesize\color{mutegray} All estimates are properties of the pinned snapshot under this
unincentivized forced-choice protocol; they are not claimed as universal traits of language models.
Thank you for reading, and for any design fire you are willing to direct at it.}

\end{document}
